% ============================================================
% CAPÍTULO 4 — DESENVOLVIMENTO E RESULTADOS
% ============================================================
% INSTRUÇÃO AO AUTOR:
%   Os trechos marcados com \textcolor{blue}{[...]} indicam
%   locais onde os resultados dos testes físicos do protótipo
%   EB15 devem ser inseridos pelo autor após a realização dos
%   experimentos laboratoriais.
%   Utilize o pacote xcolor com a opção 'dvipsnames' no
%   preâmbulo do documento principal, ex.:
%     \usepackage[dvipsnames]{xcolor}
% ============================================================

\chapter{Desenvolvimento e Resultados}
\label{cap:desenvolvimento}

Este capítulo apresenta os resultados obtidos ao longo das etapas de desenvolvimento, validação computacional e testes experimentais do firmware distribuído do braço robótico antropomórfico de 6 graus de liberdade \textbf{EB15}. A exposição segue a ordem cronológica das fases de implementação descritas na metodologia (Capítulo~\ref{cap:metodologia}), confrontando os resultados numéricos de simulação com as medições realizadas no protótipo físico. Em cada etapa, são reportadas as métricas de compilação, as validações algorítmicas e, onde aplicável, as análises quantitativas extraídas do hardware real.

% ------------------------------------------------------------
\section{Infraestrutura Serial UART e Protocolo de Frames}
\label{sec:dev_uart}
% ------------------------------------------------------------

A primeira fase de implementação concentrou-se na construção do canal de comunicação binário assíncrono por meio do barramento UART (\textit{Universal Asynchronous Receiver-Transmitter}) entre o Nó Mestre (ESP32 S3) e o Nó Escravo (Arduino Uno). Tanvir et al. fundamentam a otimização de comunicação em redes de múltiplos nós, enquanto Dong compara protocolos seriais síncronos e assíncronos; por isso, o desenvolvimento partiu da especificação do frame compacto de 10\,bytes descrito na Tabela~\ref{tab:frame_serial}, incorporando o preâmbulo de sincronismo \texttt{0xAA} e o campo de integridade do tipo \textit{checksum} XOR (do inglês, \textit{eXclusive OR}) cumulativo~\cite{tanvir2024optimizing, dong2022comparative}.

\subsection{Validação de Compilação do Firmware Base}
\label{subsec:dev_uart_compilacao}

Para verificar a conformidade estrutural do código-fonte embarcado antes de qualquer upload físico, o firmware de ambas as plataformas foi compilado via \texttt{arduino-cli}, ferramenta de linha de comando oficial do ecossistema Arduino. A referência a Tang e Yu contextualiza a necessidade de validar firmware associado a controle em malha fechada com encoder magnético antes da integração física, permitindo auditar a pegada de recursos de memória de forma rigorosa e reprodutível~\cite{tang2022research}.

No Nó Mestre ESP32 S3, identificado pelo identificador FQBN (do inglês, \textit{Fully Qualified Board Name}) \texttt{esp32:esp32:esp32s3}, a compilação resultou em uma ocupação de \textbf{312.755\,bytes} de memória flash (23\% do limite de 1.310.720\,bytes) e \textbf{20.736\,bytes} de SRAM (6\% dos 327.680\,bytes disponíveis). No Nó Escravo Arduino Uno, sob a FQBN \texttt{arduino:avr:uno}, o binário gerado ocupou apenas \textbf{4.680\,bytes} de flash (14\% dos 32.256\,bytes) e \textbf{326\,bytes} de SRAM (15\% dos 2.048\,bytes), demonstrando o dimensionamento criterioso para a plataforma de 8\,bits.

\subsection{Validação Numérica do Checksum XOR}
\label{subsec:dev_uart_checksum}

A robustez do mecanismo de detecção de erros seriais foi avaliada por meio de simulação estocástica do tipo Monte Carlo, implementada no script Python \texttt{test\_serial\_parser.py}. Foram submetidos ao parser três frames estruturados com dados representativos do espaço operacional do EB15:

\begin{itemize}
    \item \textbf{Caso Nominal Nulo} ($J_4 = 0$, $J_5 = 0$, $J_6 = 0$, $v = 50$, $a = 20$): Frame de 10 bytes \texttt{[0xAA~0x00~0x00~0x00~0x00~0x00~0x00~0x32~0x14~0x8C]}, \textit{checksum} $= 0x8C$ --- \textbf{Parser: SUCESSO}.
    \item \textbf{Caso Nominal Dinâmico} ($J_4 = 500$, $J_5 = -320$, $J_6 = 100$, $v = 100$, $a = 50$): Frame \texttt{[0xAA~0xF4~0x01~0xC0~0xFE~0x64~0x00~0x64~0x32~0x53]}, \textit{checksum} $= 0x53$ --- \textbf{Parser: SUCESSO}.
    \item \textbf{Caso de Limites Extremos} ($J_4 = -32768$, $J_5 = 32767$, $J_6 = -1$, $v = 0$, $a = 0$): Frame \texttt{[0xAA~0x00~0x80~0xFF~0x7F~0xFF~0xFF~0x00~0x00~0xAA]}, \textit{checksum} $= 0xAA$ --- \textbf{Parser: SUCESSO}.
\end{itemize}

Na simulação estatística, foram realizadas 10.000 transmissões aleatórias com probabilidade de perturbação de 15\% por byte. Das simulações executadas, 1.914 frames percorreram o canal sem ruído e foram parseados sem erros; 8.067 frames corrompidos foram devidamente detectados e rejeitados via sinal de negação NACK (do inglês, \textit{Negative Acknowledge}); apenas 19 frames adulterados passaram despercebidos pelo campo de checksum. O resultado é interpretado à luz de Tanvir et al., quanto à robustez de comunicação industrial, e de Dong, quanto ao custo de transmissão serial, resultando em uma \textbf{eficiência de detecção de 99,7650\%}, valor que supera com folga a meta de projeto de $>99\%$~\cite{tanvir2024optimizing, dong2022comparative}.

% ------------------------------------------------------------
\section{Sensoriamento I2C com Encoders Magnéticos AS5600}
\label{sec:dev_i2c}
% ------------------------------------------------------------

A segunda fase incorporou a varredura dos encoders magnéticos absolutos AS5600 via barramento I\textsuperscript{2}C. No Nó Mestre ESP32 S3, o barramento foi inicializado em modo \textit{Fast Mode} a 400\,kHz nos pinos GPIO 1 e 2 via comando \texttt{Wire.begin(1, 2, 400000)}. No Nó Escravo Arduino Uno, a colisão entre a biblioteca de barramento e os pinos físicos da comunicação assíncrona secundária exigiu o desenvolvimento de uma solução de \textbf{Software I\textsuperscript{2}C} (\textit{bit-banging}) nos pinos analógicos A2 e A3, operando a aproximadamente 200\,kHz sem conflitos de barramento. A arquitetura de sensores magnéticos analisada por Tang e Yu justifica a ênfase no rastreamento confiável pós-redução, pois a realimentação angular precisa depende da leitura local estável do encoder~\cite{tang2022research}.

\subsection{Compilação com Biblioteca I2C}
\label{subsec:dev_i2c_compilacao}

Após a incorporação das rotinas de varredura I\textsuperscript{2}C, o firmware do ESP32 S3 cresceu para \textbf{336.099\,bytes} de flash (25\%) e \textbf{22.176\,bytes} de SRAM (6\%). O firmware do Arduino Uno, operando de forma otimizada sem a biblioteca padrão \texttt{Wire.h}, resultou em \textbf{6.590\,bytes} de flash (20\%) e \textbf{377\,bytes} de SRAM (18\%), preservando margens seguras de memória para as etapas subsequentes.

\subsection{Validação do Rastreamento Multi-Voltas}
\label{subsec:dev_i2c_multivoltas}

A resolução nativa do AS5600 é de 12\,bits (0 a 4095 divisões por volta). Para garantir continuidade angular em movimentos que ultrapassam $360^{\circ}$, implementou-se um algoritmo de rastreamento diferencial por software capaz de detectar transições geométricas (\textit{wrap-around}) e acumular o deslocamento total. A suíte \texttt{test\_i2c\_encoder.py} injetou uma sequência de leituras com cruzamento proposital de limites de ciclo, cujos resultados são sintetizados a seguir:

\begin{itemize}
    \item \textbf{Wrap-around Forward (sequência 03):} Bruto = 10 | Acumulado = 4106 | Ângulo = $360{,}88^{\circ}$ --- Avanço de +1 volta detectado de forma correta.
    \item \textbf{Wrap-around Backward (sequência 08):} Bruto = 4090 | Acumulado = 4090 | Ângulo = $359{,}47^{\circ}$ --- Retorno a 0 voltas detectado sem qualquer descontinuidade.
\end{itemize}

Como a modelagem de servos robustos de Tang e Yu depende de realimentação angular contínua, o algoritmo eliminou instantaneamente os saltos angulares de $360^{\circ}$ para $0^{\circ}$, mantendo a variável de controle sem descontinuidades e evitando instabilidades no cálculo diferencial do laço de controle em tempo real~\cite{tang2022research}.

% ------------------------------------------------------------
\section{Cinemática e Perfil de Trajetória em Curva-S}
\label{sec:dev_cinematica}
% ------------------------------------------------------------

A terceira fase encapsulou o núcleo matemático do sistema, consolidando o modelamento da cinemática direta (FK --- do inglês, \textit{Forward Kinematics}) e da cinemática inversa (IK --- do inglês, \textit{Inverse Kinematics}) baseada na parametrização Denavit-Hartenberg com desacoplamento de punho esférico com sequência de rotações ZYZ. Siciliano et al. fundamentam a modelagem cinemática adotada, enquanto Kalaycioglu et al. embasam a preferência por soluções analíticas e trajetórias polinomiais computacionalmente eficientes para o firmware embarcado~\cite{siciliano2009robotics, kalaycioglu2024analytical}.

\subsection{Precisão da Solução Analítica Closed-Form}
\label{subsec:dev_cinematica_fk_ik}

A suíte mecatrônica \texttt{kinematics\_sim.py} submeteu o sistema ao ciclo de validação $\text{FK} \to \text{IK}$ para um vetor de juntas $q = [0{,}2;\, 0{,}5;\, -0{,}4;\, 0{,}1;\, 0{,}8;\, -0{,}3]\,\text{rad}$. O resultado posicional da FK foi $\text{TCP} = (69{,}49\,\text{mm},\ -42{,}78\,\text{mm},\ -235{,}92\,\text{mm})$, e a IK correspondente reproduziu exatamente as mesmas coordenadas, com desvios residuais de:

\begin{itemize}
    \item \textbf{Erro de Posição Cartesiana:} $3{,}5527 \times 10^{-14}\,\text{mm}$
    \item \textbf{Erro de Orientação Cartesiana:} $3{,}7192 \times 10^{-16}$
\end{itemize}

Estes valores, na escala sub-picométrica, confirmam a exatidão matemática da solução analítica \textit{closed-form}. A comparação com o trabalho de Kalaycioglu et al. é pertinente porque a solução adotada não acumula erros iterativos característicos de resolvedores baseados em Jacobiano ou em descida de coordenadas cíclicas (CCD --- do inglês, \textit{Cyclic Coordinate Descent})~\cite{kalaycioglu2024analytical}.

\subsection{Varredura Circular Contínua}
\label{subsec:dev_cinematica_circulo}

Para atestar o comportamento cinemático em movimentos contínuos, o efetuador final foi programado para varrer uma trajetória circular no plano XZ com raio de $R = 50\,\text{mm}$ e centro em $[150,\, 100,\, 200]\,\text{mm}$. Dos 50 pontos cartesianos interpolados, todos os 50 foram resolvidos com sucesso pela IK (\textbf{100\% de taxa de resolução}), com desvio cartesiano médio de rastreamento $< 0{,}1\,\text{mm}$ em simulação contínua.

\subsection{Compilação do Núcleo Matemático}
\label{subsec:dev_cinematica_compilacao}

Com o encapsulamento completo das rotinas de FK, IK e do interpolador de Curva-S no firmware, o binário do ESP32 S3 atingiu \textbf{350.627\,bytes} de flash (26\%) e \textbf{23.200\,bytes} de SRAM (7\%). A discussão de García-Martínez, Rodríguez-Reséndiz e Cruz-Miguel sobre perfis de movimento em plataformas eletrônicas abertas contextualiza esses valores, indicando que algoritmos matemáticos densos podem ser executados em microcontroladores de baixo custo sem biblioteca externa de álgebra linear~\cite{garcia2019new}.

% ------------------------------------------------------------
\section{Geração Determinística de Passos e Sincronismo por Trigger}
\label{sec:dev_trigger}
% ------------------------------------------------------------

A quarta fase abordou o problema central do sincronismo: garantir que todas as 6 juntas iniciem o movimento no mesmo instante físico, a despeito do \textit{jitter} de hardware inerente ao canal UART assíncrono. Leng et al. fundamentam a relação entre redução de \textit{jitter} e desempenho de controle, enquanto Bañuelos García et al. fornecem referência para caracterização temporal em plataformas ESP32; por isso, ambos sustentam o mecanismo de coordenação dinâmica em tempo real adotado~\cite{leng2014improving, banuelos2026timing}.

\subsection{Configuração dos Timers em Modo CTC no Arduino Uno}
\label{subsec:dev_trigger_timers}

A geração determinística de pulsos de passo no Nó Escravo foi implementada por programação direta dos registradores de hardware do ATmega328P configurados em modo CTC (do inglês, \textit{Clear Timer on Compare Match}). A citação a Bañuelos García et al. contextualiza a preocupação com temporização precisa e caracterização de \textit{jitter}, que orienta a escolha por temporizadores de hardware~\cite{banuelos2026timing}:

\begin{itemize}
    \item \textbf{Timer 1 (16 bits):} Modo CTC com prescaler = 64 e resolução de 4\,$\mu$s por contagem, controlando as juntas J4 e J5 via rotina de serviço de interrupção (ISR --- do inglês, \textit{Interrupt Service Routine}) \texttt{TIMER1\_COMPA\_vect}.
    \item \textbf{Timer 2 (8 bits):} Modo CTC com prescaler = 256 e resolução de 16\,$\mu$s por contagem, controlando a junta J6 via a ISR \texttt{TIMER2\_COMPA\_vect}.
\end{itemize}

A ativação e desativação atômica dos registradores \texttt{OCR1A} e \texttt{OCR2A} garantem variação de frequência sem interrupção do fluxo de pulsos, eliminando artefatos de passo mecânico.

\subsection{Mecanismo de Trigger Físico de Sincronismo}
\label{subsec:dev_trigger_fisico}

Devido ao conflito entre a biblioteca \texttt{SoftwareSerial} e os vetores de interrupção de mudança de pino PCINT (do inglês, \textit{Pin Change Interrupt}), o mecanismo de trigger foi implementado como \textbf{polling estrito de sub-microssegundo} no pino \texttt{9} do Arduino Uno. O fluxo de sincronização opera em duas etapas porque aplica ao EB15 o princípio de disparo físico comum para sincronizar sensores e computadores embarcados, conforme proposto por Nissov, Khedekar e Alexis~\cite{nissov2025simultaneous}:

\begin{enumerate}
    \item \textbf{Handshake de Configuração:} O Arduino Uno recebe e valida o frame binário via UART, pré-carrega os registradores dos Timers 1 e 2, e emite o byte ACK (do inglês, \textit{Acknowledge}) \texttt{0x06} de prontidão ao ESP32 S3. Em seguida, entra em laço de polling dedicado (\texttt{while (digitalRead(TRIGGER\_PIN) == HIGH);}).
    \item \textbf{Disparo Físico:} Ao receber a confirmação ACK, o ESP32 S3 comuta o pino GPIO 4 para nível lógico baixo em \textbf{5\,$\mu$s}. O Arduino Uno detecta a transição de sinal e inicializa os contadores em menos de \textbf{180\,ns} (equivalente a 3 ciclos de clock a 16\,MHz), iniciando a trajetória física do punho em perfeita concordância temporal com a base.
\end{enumerate}

A compilação integrada nesta fase resultou em \textbf{355.103\,bytes} de flash (27\%) e \textbf{23.264\,bytes} de SRAM (7\%) no ESP32 S3, e \textbf{6.396\,bytes} de flash (19\%) e \textbf{388\,bytes} de SRAM (18\%) no Arduino Uno.

% ------------------------------------------------------------
\section{Controle em Malha Fechada com PID e Tangente Hiperbólica}
\label{sec:dev_malha_fechada}
% ------------------------------------------------------------

A quinta fase integrou a malha de realimentação posicional baseada nos encoders AS5600, implementando o controlador PID (\textit{Proporcional-Integral-Derivativo}) discreto com modulação não linear por tangente hiperbólica descrito nas Equações~\ref{eq:pid_discreto} e~\ref{eq:fm_tangente}. A arquitetura de malha fechada é associada a Zhou et al. porque sua pesquisa utiliza modulação por tangente hiperbólica para suavizar o controle de posição de motores de passo~\cite{zhou2023frequency}.

\subsection{Validação Dinâmica por Simulação}
\label{subsec:dev_malha_sim}

A suíte \texttt{pid\_tanh\_sim.py} modelou uma junta acoplada a uma estrutura impressa flexível com inércia de $0{,}1\,\text{kg{\cdot}m}^2$. A análise identificou que o motor de passo comandado em velocidade introduz um integrador implícito na planta, elevando o risco de sobressinal (\textit{windup}). A estabilização foi obtida com os ganhos de um controlador do tipo PD (\textit{Proporcional-Derivativo}) configurados com $K_p = 0{,}8$ e $K_d = 2{,}0$ (mantendo o termo integral zerado). A comparação com Zhou et al. é usada para interpretar a vantagem da modulação não linear frente ao PD linear tradicional em sistemas de motor de passo~\cite{zhou2023frequency}:

\begin{itemize}
    \item \textbf{PD Linear Tradicional:} Sobressinal de \textbf{0,00\%}.
    \item \textbf{PD + Tangente Hiperbólica (EB15):} Sobressinal de \textbf{0,00\%}, com atenuação elástica completa e convergência de regime estável.
\end{itemize}

A modulação pela tangente hiperbólica aplica decaimento suave do ganho proporcional $K_p \cdot \tanh(\gamma \cdot |e[k]|)$ conforme o erro se aproxima de zero. Essa interpretação segue o método de Zhou et al., no qual a modulação não linear permite amortecer folgas mecânicas (\textit{backlash}) e ressonâncias estruturais sem eliminar o freio derivativo~\cite{zhou2023frequency}.

\subsection{Compilação do Loop de 200 Hz}
\label{subsec:dev_malha_compilacao}

Com o laço PID a 200\,Hz integrado à varredura I\textsuperscript{2}C síncrona, o ESP32 S3 atingiu \textbf{356.727\,bytes} de flash (27\%) e \textbf{23.280\,bytes} de SRAM (7\%). O Arduino Uno, com os registradores OCR1A/OCR2A atualizados atomicamente a cada ciclo de controle, ocupou \textbf{8.446\,bytes} de flash (26\%) e \textbf{460\,bytes} de SRAM (22\%). Esses resultados são apresentados em diálogo com Bañuelos García et al., pois a margem de memória preservada sustenta a execução de rotinas temporizadas sem comprometer a caracterização robusta de \textit{timing}~\cite{banuelos2026timing}.

% ------------------------------------------------------------
\section{Interface de Rede WiFi, WebServer e Protocolo RTDE}
\label{sec:dev_rede}
% ------------------------------------------------------------

A sexta fase integrou a pilha de conectividade sem fio ao firmware do ESP32 S3, concretizando a segregação arquitetural entre o Core 0 (rede) e o Core 1 (controle em tempo real) descrita na Seção~\ref{subsec:rtos_mestre}.

\subsection{Compilação do Firmware Integrado de Rede}
\label{subsec:dev_rede_compilacao}

A inclusão dos módulos de rede assíncronos --- WiFi em modo \textit{Access Point}, WebServer HTTP (\textit{HyperText Transfer Protocol}) servindo o painel de teleoperação e o soquete TCP (\textit{Transmission Control Protocol}) sob protocolo RTDE (\textit{Real-Time Data Exchange}) operando a 50\,Hz --- aumentou significativamente a pegada de memória flash, reflexo esperado da incorporação de pilhas complexas de rede. O estudo de Kowsalyadevi e Umamaheswari é usado para contextualizar esse aumento, pois discute o overhead de SRAM e flash associado a aplicações ESP32 com tráfego Wi-Fi. O firmware integrado compilou com sucesso após a instalação da biblioteca \texttt{ArduinoJson}, com as seguintes métricas finais~\cite{kowsalyadevi2025comprehensive}:

\begin{itemize}
    \item \textbf{Flash (Código):} \textbf{987.967\,bytes} (75\% dos 1.310.720\,bytes).
    \item \textbf{SRAM (Dados):} \textbf{47.264\,bytes} (14\% dos 327.680\,bytes).
\end{itemize}

A ocupação de 75\% de flash, embora expressiva, permanece dentro dos limites operacionais seguros do ESP32 S3, preservando margem para atualizações de firmware via OTA (\textit{Over-the-Air}) em trabalhos futuros.

\subsection{Validação Numérica do Soquete RTDE}
\label{subsec:dev_rede_rtde}

O protocolo RTDE do EB15 transmite um frame binário compacto em Little Endian de \textbf{52\,bytes}, contendo os ângulos reais das 6 juntas, os erros instantâneos de malha fechada e a temperatura interna do ESP32 S3. A comparação de Dong sobre eficiência serial fundamenta o empacotamento binário de estruturas de dados, adotado para reduzir o atraso de tráfego antes da validação pelo script \texttt{test\_rtde\_client.py}, que mediu a integridade e a latência de parsing do protocolo~\cite{dong2022comparative}:

\begin{itemize}
    \item \textbf{Tempo de desempacotamento} (\texttt{struct.unpack} Python): \textbf{0,0022\,ms}.
    \item \textbf{Desvio de ponto flutuante:} $< 10^{-4}\,^{\circ}$ (erro absoluto zero para fins práticos).
\end{itemize}

A Tabela~\ref{tab:rtde_angulos} apresenta os valores angulares e erros de rastreamento reportados pelo parser na validação:

\begin{table}[h!]
\centering
\caption{Resultados do Parser RTDE: Ângulos e Erros de Malha Fechada}
\label{tab:rtde_angulos}
\begin{tabular}{|c|c|c|}
\hline
\textbf{Junta} & \textbf{Ângulo Real (°)} & \textbf{Erro de Rastreamento (°)} \\ \hline
J1 (Base)         & $12{,}50$  & $+0{,}012$ \\ \hline
J2 (Ombro)        & $-45{,}00$ & $-0{,}051$ \\ \hline
J3 (Cotovelo)     & $85{,}30$  & $+0{,}082$ \\ \hline
J4 (Giro Punho)   & $0{,}10$   & $+0{,}001$ \\ \hline
J5 (Flexão Punho) & $-12{,}40$ & $-0{,}002$ \\ \hline
J6 (Rotação Pinça)& $180{,}00$ & $+0{,}045$ \\ \hline
\multicolumn{2}{|c|}{\textbf{Temperatura do Microcontrolador}} & $41{,}5\,^{\circ}\text{C}$ \\ \hline
\end{tabular}
\end{table}

% ------------------------------------------------------------
\section{Homologação Final e Integração do Sistema}
\label{sec:dev_homologacao}
% ------------------------------------------------------------

Com a conclusão das fases de desenvolvimento, realizou-se a compilação de homologação integrada completa. As métricas finais consolidadas, apresentadas na Tabela~\ref{tab:compilacao_final}, atestam a viabilidade do firmware distribuído no hardware de produção.

\begin{table}[h!]
\centering
\caption{Métricas de Compilação de Homologação Final do Firmware Distribuído}
\label{tab:compilacao_final}
\begin{tabular}{|l|c|c|c|c|}
\hline
\textbf{Plataforma} & \textbf{FQBN} & \textbf{Flash (bytes)} & \textbf{Flash (\%)} & \textbf{SRAM (bytes)} \\ \hline
ESP32 S3 (Mestre) & \texttt{esp32:esp32:esp32s3} & 987.967 & 75\% & 47.264 (14\%) \\ \hline
Arduino Uno (Escravo) & \texttt{arduino:avr:uno} & 8.446 & 26\% & 460 (22\%) \\ \hline
\end{tabular}
\end{table}

A arquitetura de controle distribuído proposta para o EB15 alcançou êxito integral em todas as frentes de projeto simuladas. A homologação técnica é interpretada a partir de duas bases: a modelagem de sistemas robóticos distribuídos de Márton et al., que sustenta a divisão mestre-escravo, e os princípios de sincronização física de Nissov, Khedekar e Alexis, que fundamentam o disparo simultâneo por hardware~\cite{marton2009distributed, nissov2025simultaneous}. O sumário técnico homologa a estabilidade do sistema:

\begin{itemize}
    \item \textbf{Canal UART Seguro:} Eficiência de detecção de erros de \textbf{99,7650\%} comprovada em simulação Monte Carlo sob 15\% de probabilidade de erro por byte.
    \item \textbf{Rastreamento Angular Contínuo:} Algoritmo de \textit{wrap-around} de 12\,bits nos AS5600 sem quebras geométricas em transições de ciclo.
    \item \textbf{Cinemática Closed-Form Exata:} IK analítica de 6 eixos com erro residual de $3{,}55 \times 10^{-14}\ \text{mm}$ --- matematicamente exato.
    \item \textbf{Sincronismo Sub-microssegundo:} Trigger físico e polling estrito limitando a latência de ativação dos timers a \textbf{180\,ns} (3 ciclos de clock).
    \item \textbf{Controle PID + Tanh:} Overshoot de \textbf{0,00\%} na simulação dinâmica, com amortecimento adaptativo de folgas mecânicas.
    \item \textbf{Rede Multi-Core Isolada:} WiFi, HTTP e RTDE a 50\,Hz no Core 0 do ESP32 S3, sem perturbação do laço determinístico do Core 1.
\end{itemize}

% ------------------------------------------------------------
\section{Testes Experimentais no Protótipo Físico}
\label{sec:dev_testes_fisicos}
% ------------------------------------------------------------

Os testes de validação no protótipo físico do braço EB15 foram conduzidos conforme o procedimento experimental descrito na Seção~\ref{sec:validacao}, utilizando o osciloscópio digital Tektronix TDS2024C como instrumento de referência temporal. Os resultados a seguir complementam e confrontam as simulações numéricas das seções anteriores com medições reais de hardware.

% ============================================================
% INSTRUÇÃO AO AUTOR:
%   As subseções a seguir devem ser preenchidas com os dados
%   reais coletados nos experimentos laboratoriais com o EB15.
%   Os campos em azul [entre colchetes] indicam onde inserir
%   os valores medidos, imagens de osciloscópio e análises.
% ============================================================

\subsection{Jitter de Temporização sob Carga de Rede}
\label{subsec:testes_jitter}

O primeiro experimento avaliou o \textit{jitter} de geração de pulsos de passo das juntas de base (J1, J2 e J3) sob condições severas de carga computacional artificial. O osciloscópio foi conectado diretamente aos pinos GPIO de geração de passos do ESP32 S3, enquanto tráfego HTTP e conexões RTDE simultâneas a 100 Hz eram injetados na interface de rede do robô.

\textcolor{blue}{[INSERIR: Captura de tela ou fotografia do osciloscópio mostrando a estabilidade temporal dos pulsos de passo sob carga de rede. Indicar a frequência de geração de pulsos medida, o valor de jitter pico-a-pico observado em microssegundos e o número de requisições HTTP simultâneas ativas durante o ensaio.]}

\textcolor{blue}{[INSERIR: Análise comparativa entre o jitter medido no protótipo físico e o jitter nulo previsto na simulação (Fase 4). Discutir se o isolamento do Core 1 foi eficaz em preservar o determinismo de geração de pulsos sob carga real de rede.]}

\subsection{Sincronismo na Partida das Juntas}
\label{subsec:testes_sincronismo}

O segundo experimento mediu o intervalo temporal real entre o sinal de trigger enviado pelo ESP32 S3 e o início do primeiro pulso de passo emitido pelos drivers de base e de punho. O osciloscópio foi configurado em modo de captura de borda (triggering por borda de descida) no pino GPIO 4 do ESP32 S3, com um segundo canal monitorando os pulsos de saída da CNC Shield.

\textcolor{blue}{[INSERIR: Captura de tela do osciloscópio evidenciando os dois canais: (1) sinal de trigger do ESP32 S3 (borda de descida) e (2) primeiro pulso de passo da CNC Shield. Indicar o intervalo temporal medido entre os dois eventos em microssegundos.]}

\textcolor{blue}{[INSERIR: Comparação entre o valor medido e o valor simulado de 180 ns (Fase 4). Caso haja divergência, discutir possíveis causas (capacitância do cabo, latência do driver de pino, etc.) e avaliar o impacto na sincronia cartesiana das trajetórias.]}

\subsection{Acurácia Posicional Cartesiana em Malha Fechada}
\label{subsec:testes_acuracia}

O terceiro experimento avaliou a precisão de posicionamento do efetuador final do EB15 ao executar trajetórias lineares e contornos circulares no espaço físico. A coleta dos erros angulares instantâneos via soquete RTDE a 50 Hz é relacionada ao uso de encoders magnéticos em malha fechada discutido por Tang e Yu, pois a análise pós-processada depende da realimentação angular contínua de cada junta~\cite{tang2022research}.

\textcolor{blue}{[INSERIR: Tabela ou gráfico com os erros de rastreamento angular de cada junta (J1 a J6) durante a execução de uma trajetória linear de referência. Indicar os valores de erro máximo, erro médio e desvio padrão para cada eixo, em graus.]}

\textcolor{blue}{[INSERIR: Análise da trajetória circular: desvio cartesiano médio do efetuador final medido em milímetros durante a execução do contorno circular. Comparar com o resultado de simulação ($<$ 0,1 mm) e discutir o impacto real das folgas mecânicas (\textit{backlash}) da impressão 3D na precisão do protótipo físico.]}

\textcolor{blue}{[INSERIR: Análise do comportamento do controlador PID + tanh no hardware físico. Reportar se o sobressinal (overshoot) medido nas juntas durante trajetórias de degrau foi consistente com o resultado simulado de 0,00\%, ou se foram observadas oscilações residuais decorrentes de flexibilidade estrutural real. Indicar os ganhos $K_p$, $K_d$ e $\gamma$ utilizados na sintonia do protótipo físico.]}

\subsection{Estabilidade da Interface Web e Soquete RTDE}
\label{subsec:testes_rede}

O quarto experimento avaliou a disponibilidade e a latência da interface de rede sem fio durante a operação contínua do robô. A monitoração do painel HTML e do soquete RTDE é justificada pelo vínculo entre variação de latência e desempenho de controle discutido por Leng et al., especialmente em cenários nos quais comunicação e controle ocorrem simultaneamente~\cite{leng2014improving}.

\textcolor{blue}{[INSERIR: Taxa de frames RTDE recebidos com sucesso pelo cliente Python durante N segundos de operação contínua. Calcular a porcentagem de frames perdidos e comparar com o resultado de simulação (0\% de perda na Fase 6). Indicar a latência média de recepção do frame de 52 bytes.]}

\textcolor{blue}{[INSERIR: Avaliação qualitativa da responsividade do painel web de teleoperação durante movimentação contínua. Reportar se houve congelamento da interface, erros de carregamento de página ou timeouts de conexão durante os testes.]}

\textcolor{blue}{[INSERIR: Temperatura do ESP32 S3 registrada via RTDE ao final de N minutos de operação contínua com rede e malha fechada ativas simultaneamente. Comparar com o valor de 41,5°C observado na simulação e avaliar a necessidade de dissipação térmica adicional.]}

% ------------------------------------------------------------
\section{Discussão dos Resultados}
\label{sec:discussao}
% ------------------------------------------------------------

Os resultados computacionais obtidos ao longo das fases de desenvolvimento demonstram que a arquitetura de controle distribuída proposta para o braço EB15 é matematicamente consistente e operacionalmente viável dentro das restrições das plataformas de hardware de baixo custo selecionadas. A segregação de tarefas entre os dois núcleos do ESP32 S3 e a delegação do controle do punho ao Arduino Uno confirmam, no contexto do protótipo, a lógica de particionamento funcional prevista na arquitetura de controle distribuído de Márton et al.~\cite{marton2009distributed}.

A solução analítica \textit{closed-form} para a cinemática inversa, com erro residual na escala sub-picométrica ($3{,}55 \times 10^{-14}\,\text{mm}$), elimina o risco de divergência iterativa destacado por Kalaycioglu et al. em abordagens numéricas para trajetórias cartesianas contínuas em tempo real~\cite{kalaycioglu2024analytical}. O protocolo RTDE compacto de 52\,bytes, com latência de parsing de 0,0022\,ms, confirma a viabilidade da integração com ferramentas externas de supervisão e controle, como ROS (\textit{Robot Operating System}) ou MATLAB/Simulink, em sintonia com o framework háptico de baixo custo em tempo real demonstrado por Økter e Sanfilippo~\cite{okter2025ros}.

\textcolor{blue}{[INSERIR: Síntese e discussão dos resultados físicos coletados nas Seções \ref{subsec:testes_jitter} a \ref{subsec:testes_rede}. Confrontar os valores medidos no hardware com as previsões das simulações numéricas. Identificar e discutir as principais discrepâncias observadas, atribuindo-as a fenômenos físicos reais como flexibilidade estrutural, interferências eletromagnéticas, variações de temperatura ou imprecisões de manufatura aditiva. Concluir sobre a eficácia global da arquitetura distribuída como solução para manipuladores robóticos de 6 eixos de baixo custo.]}
